\section*{Capítulo \ref{ch:b1:fluid-envelopes-structure} --- Estrutura da atmosfera e do oceano}
% values: calc/ch26.py

\begin{solution}{exo:b1:fluid-envelopes-structure:1}
(a)~$H = R^*T/(M_0 g_0) = 8314.32 \times 250/(28.9644 \times 9.80665) =
\qty{7.32}{km}$. (b)~$e^{10/7.32} = 3.92$: a pressão cai para cerca de um quarto.
(c)~$H$ é proporcional à temperatura: a \qty{216.65}{K} ela vale
\qty{6.34}{km}, contra \qty{8.43}{km} a \qty{288.15}{K}. Ar frio é mais denso a uma
dada pressão, então a pressão cai mais depressa com a altura.
\end{solution}

\begin{solution}{exo:b1:fluid-envelopes-structure:2}
Com $\Gamma_d = \qty{9.76}{K/km}$: \qtyrange{0}{0.3}{km}, $11.0 > 9.76$,
instável; \qtyrange{0.3}{1.2}{km}, $9.7$, neutra dentro da precisão dos dados;
\qtyrange{1.2}{1.5}{km}, uma inversão; \qtyrange{1.5}{8}{km}, estável; acima de \qty{8}{km},
isotérmica, estável. A inversão é a mais estável, já que $\Gamma_d - \Gamma = 13.8$ é ali o maior.
\end{solution}

\begin{solution}{exo:b1:fluid-envelopes-structure:3}
(a)~$\sigma = 1026.95 - 1000 = \qty{26.95}{kg.m^{-3}}$. (b)~B, a mais leve, fica acima de A.
(c)~Todas as densidades da água do mar começam com os mesmos dígitos, 10; $\sigma$ guarda apenas
os dígitos que mudam de uma água para outra.
\end{solution}

\begin{solution}{exo:b1:fluid-envelopes-structure:4}
(a)~$S = 0.25 \times 35.40 + 0.75 \times 34.30 = 34.575$; $\theta = 0.25 \times
15.0 + 0.75 \times 4.0 = \qty{6.75}{\celsius}$. (b)~No segmento PQ, a um quarto
do caminho de Q até P (\cref{prop:b1:fluid-envelopes-structure:mixing}).
(c)~As misturas de P, Q e R preenchem o triângulo PQR. Escrevendo
$(34.60, 8.0) = a\,\mathrm P + b\,\mathrm Q + c\,\mathrm R$ com $a + b + c = 1$
dá $a = 0.345$, $b = 0.788$, $c = -0.133$: uma proporção negativa de R, de modo que o
ponto fica fora do triângulo e nenhuma mistura das três pode produzi-lo.
\end{solution}

\begin{solution}{exo:b1:fluid-envelopes-structure:5}
(a)~$\theta = 253.15 \times 2^{2/7} = \qty{308.6}{K}$. (b)~A \qty{1000}{hPa},
$\theta = T = \qty{298.15}{K}$. A \omterm{def:b1:fluid-envelopes-structure:potential-temperature}{temperatura potencial} aumenta para cima, de
298.2 a \qty{308.6}{K}: a camada é estável para ar seco. (c)~Ele chegaria à
sua \omterm{def:b1:fluid-envelopes-structure:potential-temperature}{temperatura potencial}, \qty{308.6}{K}, isto é \qty{35.4}{\celsius}, dez
graus mais quente que o ar que já está ali.
\end{solution}

\begin{solution}{exo:b1:fluid-envelopes-structure:6}
(a)~$T = 288.15 - 6.5 \times 5 = \qty{255.65}{K}$ e $p = \num{101325} \times
(255.65/288.15)^{5.256} = \qty{54020}{Pa}$, \qty{540.2}{hPa}. (b)~$H =
\qty{8.43}{km}$ e $p = \num{101325}\,e^{-5/8.43} = \qty{56010}{Pa}$. A pressão real
(padrão) é menor: o ar ali é mais frio que \qty{288.15}{K}, logo mais denso, e a pressão
cai mais depressa com a altura. (c)~$p = \num{22632}\,
e^{-9/6.34} = \qty{5475}{Pa}$, o valor de
\cref{tab:b1:fluid-envelopes-structure:layers}.
\end{solution}

\begin{solution}{exo:b1:fluid-envelopes-structure:7}
(a)~$\beta\,\Delta S = \alpha\,\Delta T$ dá $\Delta S = 1.67/7.61 = 0.22$.
(b)~$\Delta\rho = 1026.95 \times (-1.67\times10^{-4} \times 2 + 7.61\times10^{-4}
\times 0.3) = \qty{-0.11}{kg.m^{-3}}$: a água fica mais leve.
(c)~$\Delta\rho = 1028.11 \times (-0.53\times10^{-4} \times 2 + 7.85\times10^{-4}
\times 0.3) = \qty{+0.13}{kg.m^{-3}}$: as mesmas mudanças agora a tornam mais densa, porque
o aquecimento quase não muda a densidade da água perto do congelamento.
\end{solution}

\begin{solution}{exo:b1:fluid-envelopes-structure:8}
(a)~$N^2 = (9.81/260) \times (9.76 - 4)\times10^{-3} = \qty{2.17e-4}{s^{-2}}$,
$N = \qty{0.0147}{s^{-1}}$, período $2\pi/N = \qty{426}{s}$, cerca de sete
minutos. (b)~$N^2 = (9.81/270) \times 14.76\times10^{-3} = \qty{5.36e-4}{s^{-2}}$,
$N = \qty{0.0232}{s^{-1}}$, período \qty{271}{s}. (c)~$N$ tende a zero e o período
cresce sem limite: a força restauradora se anula em uma camada neutra.
\end{solution}

\begin{solution}{exo:b1:fluid-envelopes-structure:9}
(a)~É mais densa, $\sigma_\theta = 27.20$ contra 26.30: a AAIW é mais doce mas muito mais fria.
(b)~$\Delta\sigma = 1026.95 \times [-1.67\times10^{-4} \times
(4.31 - 15.43) + 7.61\times10^{-4} \times (34.30 - 35.54)] = 1.907 - 0.969 =
0.94$, contra $0.90$ exatamente. (c)~Temperatura: o resfriamento acrescenta
\qty{1.91}{kg.m^{-3}}, a diminuição de salinidade retira \qty{0.97}{kg.m^{-3}}.
\end{solution}

\begin{solution}{exo:b1:fluid-envelopes-structure:10}
(a)~Por \cref{thm:b1:fluid-envelopes-structure:stability}, $(g/\theta)\,\mathrm
d\theta/\mathrm dz = (g/T)(\Gamma_d - \Gamma)$; $\mathrm d\theta/\mathrm dz = 0$
dá $\Gamma = \Gamma_d$. (b)~$T = T_s - \Gamma_d z$, uma camada de gradiente constante
$L = \Gamma_d$: pela \cref{prop:b1:fluid-envelopes-structure:layer-pressure}, $p =
p_s (T/T_s)^{gM/(R\Gamma_d)}$, e $gM/(R\Gamma_d) = gMc_p/(Rg) = c_pM/R$
($= 3.5$ para o ar seco). (c)~$p$ se anula onde $T = 0$, em $z = T_s/\Gamma_d =
300/9.76 = \qty{30.7}{km}$. Camadas reais de \omterm{def:b1:fluid-envelopes-structure:potential-temperature}{temperatura potencial} uniforme têm
no máximo alguns quilômetros de espessura: são limitadas por ar estável, como a
troposfera é pela estratosfera.
\end{solution}

\begin{solution}{exo:b1:fluid-envelopes-structure:11}
(a)~$S = (33.787 + 36.050)/2 = 34.919$, $\theta = \qty{8.0}{\celsius}$.
(b)~Com uma equação de estado linear, $\sigma$ é uma função afim de $S$ e
$\theta$, de modo que o $\sigma$ da mistura é a média dos dois, 27.00.
(c)~$\alpha$ vale \qty{0.75e-4}{K^{-1}} a \qty{2}{\celsius} e
\qty{2.07e-4}{K^{-1}} a \qty{14}{\celsius}, enquanto $\beta$ muda pouco
(\num{7.80e-4} e \num{7.53e-4}). A densidade cai cada vez mais depressa à medida que a
temperatura sobe: é uma função côncava da temperatura. Ao longo da reta de mistura,
a mistura está na temperatura média, onde a densidade fica acima da corda que une os dois
originais, isto é acima de 27.00. (d)~O excesso vem da curvatura, um termo de segunda ordem:
ele cresce com o quadrado da diferença de temperatura dos originais.
\end{solution}

\begin{solution}{exo:b1:fluid-envelopes-structure:12}
(a)~A diferença de pressão é $5604 - 5500 = \qty{104}{dbar}$: a porção esfria
em $0.129 \times 0.104 = \qty{0.013}{\celsius}$, até \qty{0.718}{\celsius}.
(b)~A água em volta, a \qty{5400}{m}, está a \qty{0.730}{\celsius}, mais quente: a
porção, de salinidade quase igual, é mais densa que seu entorno e afunda de volta. A
camada é estável, embora a temperatura medida aumente para baixo; sua \omterm{def:b1:fluid-envelopes-structure:potential-temperature}{temperatura
potencial} diminui para baixo. (c)~$\theta \approx
0.7313 - 0.086 \times 5.604 = \qty{0.249}{\celsius}$, a menos de uma milésima de grau do valor exato.
\end{solution}

\begin{solution}{pb:b1:fluid-envelopes-structure:1}
\begin{enumerate}
  \item A temperatura é quase uniforme até \qty{50}{m}, cai depressa até \qty{300}{m}, mais
        devagar até \qty{1000}{m}, e é quase uniforme abaixo. A salinidade diminui de 35.80 na
        superfície a um mínimo de 34.32 a \qty{800}{m}, depois aumenta para 34.90 a \qty{3000}{m}.
  \item $22.40 - 0.5 = \qty{21.90}{\celsius}$ é atingido entre \qty{50}{m} (22.10) e
        \qty{75}{m} (19.20): $50 + 25 \times 0.20/2.90 = \qty{52}{m}$.
  \item $(22.10 - 13.00)/250 = \qty{0.0364}{\celsius/m}$, \qty{3.6}{\celsius} por \qty{100}{m}.
        A \omterm{def:b1:fluid-envelopes-structure:ocean-layers}{termoclina} se estende da base da \omterm{def:b1:fluid-envelopes-structure:ocean-layers}{camada de mistura}, perto de \qty{50}{m}, até cerca de
        \qty{800}{m}; é mais nítida entre \qty{50}{m} e \qty{75}{m} (\qty{11.6}{\celsius} por
        \qty{100}{m}).
  \item $\sigma = 1026.95 \times [1 - 1.67\times10^{-4} \times 12.4 +
        7.61\times10^{-4} \times 0.80] - 1000 = 25.45$, \qty{0.74}{kg.m^{-3}} alta demais. A água
        da superfície está longe do estado de referência, e $\alpha$ é muito maior a
        \qty{22}{\celsius} que a \qty{10}{\celsius}: a forma linear subestima o alívio causado
        pelo calor.
  \item O mínimo de salinidade a \qty{800}{m}, 34.32 a \qty{4.8}{\celsius}, é a Água Intermediária
        Antártica, formada sob os ventos de oeste chuvosos do Oceano Austral. O máximo a
        \qty{3000}{m}, 34.90 a \qty{2.5}{\celsius}, é a Água Profunda do Atlântico Norte, formada
        nos mares da Groenlândia e de Labrador.
  \item $(2.50 - 2.26)/3.042 = \qty{0.079}{\celsius}$ por \qty{1000}{dbar}.
  \item $\sigma_\theta$ aumenta com a profundidade: a coluna é estável. A \omterm{def:b1:fluid-envelopes-structure:ocean-layers}{picnoclina} é onde
        $\sigma_\theta$ aumenta mais depressa, nas poucas centenas de metros superiores, junto com a
        \omterm{def:b1:fluid-envelopes-structure:ocean-layers}{termoclina}: \num{0.82} entre \qty{200}{m} e \qty{800}{m}, apenas \num{0.70} sobre os
        \qty{2200}{m} abaixo.
  \item $S = (35.45 + 34.32)/2 = 34.885$; $\theta = (14.90 + 4.80)/2 =
        \qty{9.85}{\celsius}$.
  \item C: 26.461; I: 27.310; mistura: 26.886, exatamente a média dos dois (26.886), já que a
        forma linear é afim.
  \item A média dos originais é $(26.342 + 27.159)/2 = 26.751$; a mistura, a 26.888, é
        \qty{0.14}{kg.m^{-3}} mais densa. Isto é o cabbeling.
  \item Com a forma linear, a descida dá 26.61 a \qty{300}{m}, 26.88 a \qty{500}{m} e 27.31 a
        \qty{800}{m}: a densidade da mistura fica perto de \qty{500}{m}. A água ali,
        $(34.82, \qty{9.6}{\celsius})$, está muito próxima da mistura $(34.885, \qty{9.85}{\celsius})$: a descida entre a Água Central e a AAIW é ela mesma próxima de uma
        reta de mistura entre elas.
  \item $S = 0.25 \times 35.45 + 0.75 \times 34.32 = 34.60$, $\theta = 0.25
        \times 14.90 + 0.75 \times 4.80 = \qty{7.33}{\celsius}$. Interpolando a descida entre
        \qty{500}{m} e \qty{800}{m} dá $(34.57, \qty{7.2}{\celsius})$ perto de \qty{650}{m}: quase
        sobre ela.
  \item Em \si{K/km}: A--B 10.0; B--C 9.7; C--D $-5.5$; D--E 6.29; E--F 6.50;
        F--G 6.57; G--H 0.33; H--I $-0.50$; I--J $-0.83$.
  \item A--B instável; B--C neutra; C--D uma inversão; D--E, E--F, F--G estáveis; G--H quase
        isotérmica, estável; H--I e I--J inversões. A tropopausa fica em G, a \qty{12.5}{km} e
        \qty{-59.0}{\celsius}, onde a temperatura para de diminuir.
  \item A temperatura média é \qty{214.65}{K}: $H = 287 \times 214.65/9.81 =
        \qty{6.28}{km}$, e $147\,e^{-4/6.28} = \qty{77.7}{hPa}$, consistente com os 78 medidos.
  \item O ar em A, a \qty{22.0}{\celsius}, é mais frio que o mar: o calor flui do mar para o ar,
        que é aquecido por baixo. A camada A--B, com um gradiente acima de $\Gamma_d$, é o resultado.
  \item Em B: $22.0 - 9.76 \times 0.5 = \qty{17.1}{\celsius}$, contra 17.0; em C:
        \qty{7.4}{\celsius}, contra 7.3; em D: \qty{3.5}{\celsius}, contra 9.5. A porção é um
        pouco mais quente que o ar até C, e muito mais fria em D.
  \item Na base da inversão, perto de \qty{1.5}{km}: ao entrar na inversão, a porção fica mais
        fria, logo mais densa, que o ar em volta, o que interrompe sua subida.
  \item As porções liberadas pela camada superficial instável sobem através dela e a misturam até
        que sua \omterm{def:b1:fluid-envelopes-structure:potential-temperature}{temperatura potencial} seja uniforme, o que torna seu gradiente $\Gamma_d$; a
        inversão acima interrompe a mistura.
  \item Em K: A 294.15; B 293.99; C 294.08; D 300.51; E 312.65; F 329.29; G 346.28; H 369.50;
        I 446.99; J 600.91.
  \item $\theta$ diminui de A para B (instável), é quase constante de B para C (neutra) e aumenta
        acima (estável), mais depressa de C para D e acima de G: a mesma classificação da
        questão~14.
  \item E--F: $(329.29 - 312.65)/4 = \qty{4.2}{K/km}$; H--I: $(446.99 -
        369.50)/4 = \qty{19.4}{K/km}$, quase cinco vezes mais.
  \item Com a média $\theta = \qty{408.2}{K}$: $N^2 = (9.81/408.2) \times
        0.01937 = \qty{4.66e-4}{s^{-2}}$, $N = \qty{0.0216}{s^{-1}}$, um período de
        \qty{291}{s}, cerca de cinco minutos.
  \item $N^2 = (9.81/1027) \times (26.35 - 24.79)/150 = \qty{9.9e-5}{s^{-2}}$, $N = \qty{0.0100}{s^{-1}}$, um período de \qty{630}{s}. A estratosfera inferior é a mais estável,
        com $N^2$ quase cinco vezes maior que na \omterm{def:b1:fluid-envelopes-structure:ocean-layers}{termoclina}.
  \item Em G, $\theta = 214.15 \times (1000/186)^{2/7} = \qty{346}{K}$: a \omterm{def:b1:fluid-envelopes-structure:potential-temperature}{temperatura potencial} na
        tropopausa é de cerca de \qty{346}{K}.
\end{enumerate}
\end{solution}
